\section{Questão 3}
Seja $X$ uma variável aleatória com PDF dada por
%
\begin{equation*}
    f_X(x) = 
    \begin{cases}
        2 \ln(a) x^2, & |x| \leq 1 \\
        0,            & |x| > 1
    \end{cases}
\end{equation*}

\begin{itemize}
    \item [a)] 
    Determine o valor da constante $a$.
    
    \item [b)] 
    Qual é a faixa dinâmica da variável randômica $X$?
    
    \item [c)] 
    Determine $\E{X}$.
    
    \item [d)] 
    Determine $\Var{X}$.
    
    \item [e)] 
    Determine $P(X > 0)$.
\end{itemize}

\respostas
\begin{itemize}
    \item [a)]
    Modula-se $a$ para satisfazer $\int_\mathbb{R} f_X(x)\ dx = 1$ (axioma da probabilidade):
    %
    \begin{align}
        \int_\mathbb{R} f_X(x) \ dx &= 1 \iff \\
        \int_{-1}^1 2\ln(a) x^2 \ dx &= 1 \iff \\
        2\ln(a) \int_{-1}^1 x^2 \ dx &= 1 \iff \\
        2\ln(a) \frac{x^3}{3}\bigg|_{-1}^1 \ dx &= 1 \iff \\
        \frac{4}{3} \ln(a) &= 1 \iff 
        a = e^{3/4}.
    \end{align}

    O resultado acima permite reescrever a PDF como
    %
    \begin{equation}
        f_X(x) = 
        \begin{cases}
            3x^2/2, & |x| \leq 1 \\
            0,      & |x| > 1
        \end{cases}
    \end{equation}
    
    \item [b)]
    A faixa dinâmica é a porção do domínio da PDF cuja imagem não é trivial: $[-1, 1]$.
    
    \item [c)]
    Sendo a PDF simétrica, $\E{X}=0$.
    
    \item [d)]
    Como $\E{X}=0$ pelo item (c), tem-se
    %
    \begin{equation}
        \Var{X} 
        = E[X^2]
        = \int_\mathbb{R} x^2 f_X(x) \ dx
        = \int_{-1}^1 x^2 \cdot \left(\frac{3}{2}x^2\right) \ dx
        = \frac{3}{2} \int_{-1}^1 x^4 \ dx
        = \frac{3}{2} \frac{x^5}{5}\bigg|_{-1}^1
        = \frac{3}{5}
    \end{equation}
    
    \item [e)]
    De maneira geral, $P(X<0) = 1 - P(X>0)$.
    Como a PDF em questão é simétrica, $P(X<0) = P(X>0)$, 
    e portanto $P(X>0) = 1/2$ ou 50\%.
\end{itemize}